\scanpage{7}

8) \deleted{\unsure{Polar|Pour}} La donnée d'une prépolarisation $P$ se décompose en des composantes $P^i$ ($i\in\mathbb Z$). Elles sont à priori tt à fait indépendantes. Une prépolarisation est dite admissible si elle satisfait les conditions

\[
\begin{aligned}
&\varphi\in P(M),\ \psi\in P(N),\ \chi\in P(M\otimes N)\\
&\quad\longrightarrow\unsure{\chi|\varphi}\ \text{compatible avec }\varphi\otimes\psi\quad\text{i.e. dans }\widehat P(M\otimes N)
\end{aligned}\tag*{(P2)}
\]

\[
\begin{aligned}
&\forall i,\quad\text{l'hom can }K(\unsure{1|i})\otimes K(\unsure{1|i})\to K(\unsure{2i|2|1i})\quad\deleted{\text{\unsure{prescrit}}}\\
&\hspace{13em}\text{est dans }\widehat P(K(\unsure{1|i})).
\end{aligned}\tag*{(P3)}
\]

En présence de (P2), (P3) équivaut à

\deleted{$(\mathrm{P}'3)_0$ \inserted{l'hom can} $K(0)\otimes K(0)\to K(0)$ est dans $\widehat P(K(0))$}

\[
\text{Si }\varphi\in P(M),\quad\text{alors }\varphi(i)\in\widehat P(M(i)).\tag*{$(\mathrm P'3)_{\unsure{3|i|1}}$}
\]

\deleted{(Il suffit de le supposer pour $i=1$) \unsure{et}}

Une \uline{polarisation} est une prépolarisation maximale et admissible.

\uline{NB} Si $P$ est une prépolarisation, elle est admissible sss $\widehat P$ est une polarisation. Ceci provient du fait suivant : si $\varphi\in\operatorname{Hom}(M\otimes M,L)$, $\varphi'\in\operatorname{Hom}(M'\otimes M',L')$ ($L,L'$ inv), et si $\varphi,\varphi',\varphi\otimes\varphi'\in\operatorname{Hom}(\underbrace{M\otimes M'}\otimes(M\otimes M'),L\otimes L')$ sont de Weil, alors $\varphi_1$ comp. à $\varphi$ et $\varphi'_1$ comp. à $\varphi'$ implique $\varphi_1\otimes\varphi'_1$ compatible à $\varphi\otimes\varphi'$.

% slanted
\margin{\uline{vérifier}}

9) \deleted{Soit} $P$ une polarisation. Alors on a, pour les formes \inserted{de polarisation}

a) stabilité par $\varphi\oplus\psi$ \inserted{a$'$) \deleted{\unsure{st}} stabilité par $\varphi+\psi$ [résulte de a) et b)]}

b) stabilité par $\varphi|N$

c) stabilité par $\varphi\otimes\psi$

e) stabilité par $\check\varphi$

d) stabilité par $\varphi(i)$

f) \deleted{Contient} $K(0)\otimes K(0)\xrightarrow{\sim}K(0)$ est de \unsure{pol} sss il est défini par $\lambda\in K$, avec $\lambda>0$

\margin{\uline{NB} les conditions b) c) e) f) impliquent déjà que $P$ est une prépolarisation. Les $P$ maximaux satisfaisant b) c) e) f) sont les prépolarisations maximales satisfaisant (P2). Donc les polarisations $\Rightarrow$ systèmes \deleted{\unsure{maximaux}} satisfaisant b) \unsure{c)} d) e) f) et \uline{maximaux} (\uline{NB} a) en résulte)}

% archive sheet number: 6
